\documentclass{article}
\usepackage[utf8]{inputenc}
\usepackage[T1]{fontenc}
\usepackage[english,ukrainian]{babel}
\usepackage{url}
\usepackage{amsmath,amssymb,amsthm}

\begin{document}

\title{Рання історія Coq}
\author{Namdak Tonpa}
\date{Травень 2025}
\maketitle

\begin{abstract}
Ця стаття присвячена ранній історії системи Coq — одного з найвпливовіших асистентів доведення у світі. Розглядаються логічні та теоретичні основи Числення індуктивних конструкцій, ключові етапи розвитку від перших прототипів до версій 7.x, а також внесок основних розробників: Тьєррі Кокана, Жерара Юе, Крістін Полен-Морінг та багатьох інших.
\end{abstract}

\section{Рання історія Coq}

\subsection{Історичні корені}

Coq — це асистент доведення для логіки вищого порядку, який дозволяє розробляти комп’ютерні програми, що відповідають їхнім формальним специфікаціям. Він є результатом приблизно десяти років досліджень проєкту Coq. Тут ми коротко розглянемо три основні аспекти: логічну мову, якою ми пишемо наші аксіоматизації та специфікації, асистент доведення, який дозволяє розробляти верифіковані математичні доведення, та екстрактор програм, який синтезує комп’ютерні програми, що відповідають їхнім формальним специфікаціям, записаним як логічні твердження в цій мові.

Логічна мова, яку використовує Coq, є різновидом теорії типів, відомою як Числення індуктивних конструкцій. Не повертаючись до Лейбніца та Буля, ми можемо датувати створення того, що зараз називається математичною логікою, працями Фреге та Пеано на межі століть. Виявлення антиномій при вільному використанні предикатів чи принципів розуміння спонукало Расселла обмежити предикатне числення стратифікацією типів. Ці зусилля завершилися створенням \emph{Principia Mathematica}, першою систематичною спробою формального обґрунтування математики. Спрощення цієї системи за принципами просто типізованого $\lambda$-числення відбулося з Простою теорією типів Чьорча.

Нотація $\lambda$-числення, спочатку використана для вираження функціональності, також могла застосовуватися як кодування доведень у природній дедукції. Цей ізоморфізм Керрі–Говарда використав Н. де Брюйн у проєкті Automath — першій повномасштабній спробі розробити та механічно верифікувати математичні доведення. Ці зусилля завершилися верифікацією Юттінгом \emph{Grundlagen} Ландау у 1970-х роках.

Використовуючи цей ізоморфізм Керрі–Говарда, значні досягнення в теорії доведень привели до появи двох типо-теоретичних систем:
\begin{itemize}
  \item Інтуїціоністська теорія типів Мартіна-Льофа, яка намагається створити нову основу математики на конструктивних принципах;
  \item Поліморфне $\lambda$-числення Жирара — дуже потужна функціональна система, в якій ми можемо представляти структури доведень логіки вищого порядку.
\end{itemize}

Поєднавши обидві системи у вищому порядку розширення мови Automath, Т. Кокан у 1985 році представив першу версію Числення конструкцій (CoC). Ця потужна логічна система дозволяла сильні аксіоматизації, але прямі індуктивні визначення були неможливі, і індуктивні поняття доводилося визначати опосередковано через функціональні кодування.

Формалізм був розширений у 1989 році Т. Коканом і К. Полен із додаванням примітивних індуктивних визначень, що призвело до сучасної Числення індуктивних конструкцій (CIC).

\subsection{Версії з 1 по 5}

\begin{note}
Цей огляд був написаний у 1995 році разом із попереднім розділом і складав початкову версію розділу «Кредити». Більш детальний опис цих ранніх версій доступний у наступних підрозділах, які взяті з документа, написаного у вересні 2015 року Жераром Юе, Тьєррі Коканом і Крістін Полен.
\end{note}

Перша реалізація CoC була розпочата у 1984 році Ж. Юе та Т. Коканом. Мовою реалізації був CAML. Ядром цієї системи був перевірник доведень для CoC, названий Конструктивним двигуном. Ця система (версія 4.10) була випущена у 1989 році. Потім система була розширена для роботи з новим численням із індуктивними типами К. Полен. Ця система (версія 5.6) була випущена у 1991 році.

Coq був перенесений на Caml-light у 1992 році (версія 5.7). Версія 5.8 вийшла у травні 1993 року. Версія V5.10 базується на загальній системі для маніпуляції термами зі зв’язуючими операторами, розробленій Четом Мурті.

Восени 1994 року К. Полен-Морінг замінила структуру індуктивно визначених типів і сімейств на нову структуру, що дозволяє взаємно рекурсивні визначення.

\begin{quote}
Роккенкур, 1 лютого 1995 року \\
Жерар Юе
\end{quote}

\subsubsection{Версія 1}

Прототип перевірника типів для Теорії конструкцій. Версія 1.10 заморожена 22 грудня 1984 року — використовувалася для дисертації Тьєррі Кокана.

\subsubsection{Версія 2}

Числення конструкцій з доведенням консистентності. Версія 2.13 з універсумами заморожена 25 червня 1986 року.

\subsubsection{Версія 3}

Початок автоматизації доведень з пошуковим алгоритмом у стилі PROLOG. Реалізація на CAML.

\subsubsection{Версія 4}

Початок витягування програм. Введення індуктивних типів (V4.4, березень 1988). Версія 4.10 — зріла система, випущена 1 травня 1989 року.

\subsubsection{Версія 5}

Перейменована в Coq для Числення індуктивних конструкцій. Значна переробка архітектури Четом Мурті (версії 5.6–5.10).

\begin{quote}
Вересень 2015 року \\
Тьєррі Кокан, Жерар Юе і Крістін Полен-Морінг
\end{quote}

\subsection{Версії 6}

\subsubsection{Версія 6.1}

Перенесення на Objective Caml. Багато нових інструментів і тактик (листопад 1996).

\subsubsection{Версія 6.2}

Перехід на camlp4, покращення синтаксису, тактик і документації (травень 1998).

\subsubsection{Версія 6.3}

Нові тактики, розширення фікс-пойнтів, покращення уніфікації (грудень 1999).

\subsection{Версії 7}

Значна ревізія архітектури (2001–2003). Нове ядро перевірки типів, модульна система, coqdoc, витягування в Caml/Haskell, тактики Field, Fourier, Omega та багато іншого.

\begin{quote}
Орсе, травень 2002 року \\
Х’юго Хербелін і Крістін Полен
\end{quote}

\end{document}
